\documentclass[a4paper,11pt]{article}

% Packages essentiels
\usepackage[utf8]{inputenc}
\usepackage[french]{babel}
\usepackage[T1]{fontenc}
\usepackage{amsmath, amssymb, amsthm}
\usepackage{array}
\usepackage{graphicx}
\usepackage{geometry}
\usepackage{xcolor}
\usepackage{fancyhdr}
\usepackage{enumitem}
\usepackage{siunitx} % Pour les unités physiques
\usepackage[european, straightvoltages, RPvoltages]{circuitikz} % Pour dessiner des circuits
\usepackage{multicol}
\usepackage{pgfplots}
\pgfplotsset{compat=1.18}


\usepackage{listings}
\usepackage{xcolor}


\definecolor{backcolour}{rgb}{0.95,0.95,0.92}
\lstset{
  language=matlab,
  backgroundcolor=\color{backcolour},   
  basicstyle=\ttfamily\small,
  keywordstyle=\color{magenta},
  stringstyle=\color{blue},
  commentstyle=\color{gray},
  numbers=left,
  numberstyle=\tiny\color{gray},
  numbersep=5pt,
  frame=single,
  breaklines=true,
  captionpos=b,
  tabsize=1,
  showspaces=false,
  showstringspaces=false
}


% Configuration de la page
\geometry{margin=2.5cm}
\pagestyle{fancy}
\fancyhf{}
\lhead{TP Programmation des Microcontroleurs}
\rhead{Préparation}

\fancyfoot{}
\lfoot{SN421 TP1}
\rfoot{Mathéo Guillevic \& Eva Berhault}
\cfoot{\thepage}

% Environnements personnalisés
\usepackage{tcolorbox}
\tcbuselibrary{theorems}
\usetikzlibrary{babel}

% Boîte pour les exercices
\newtcolorbox{exercice}[1]{
	colback=blue!5!white,
	colframe=blue!75!black,
	title=Exercice #1,
	fonttitle=\bfseries
}

% Boîte pour les solutions
\newtcolorbox{solution}{
	colback=green!5!white,
	colframe=green!50!black,
	title=Solution,
	fonttitle=\bfseries
}

% Boîte pour les rappels
\newtcolorbox{rappel}{
	colback=orange!5!white,
	colframe=orange!75!black,
	title=Rappel,
	fonttitle=\bfseries
}

% Boîte pour les infos
\newtcolorbox{info}{
	colback=blue!5!white,
	colframe=blue!75!black,
	title=Info,
	fonttitle=\bfseries
}

% Boîte pour les notes
\newtcolorbox{note}{
	colback=cyan!5!white,
	colframe=cyan!75!black,
	title=Note,
	fonttitle=\bfseries
}

% Informations du document
\title{\textbf{TP1 }\\
	\large }
\author{Mathéo Guillevic \& Eva Berhault\\4APP TP1}
\date{\today}

\begin{document}

    \begin{titlepage}
    \begin{tikzpicture}[remember picture, overlay]
        % --- Noms (Haut Gauche) ---
        \node[anchor=north west, align=left] at ([shift={(1.2cm,-1.5cm)}]current page.north west) {
            \large Eva Berhault \\
            \large Mathéo Guillevic
        };

        % --- Infos (Haut Droite) ---
        \node[anchor=north east, align=right] at ([shift={(-1.2cm,-1.5cm)}]current page.north east) {
            \large Date : 14/09/2026\\
            \large Gpe : 1
        };

        % --- Numéro de TP (Haut Centre) ---
        \node[anchor=north] at ([shift={(0cm,-2.5cm)}]current page.north) {
            \large TP n$^\circ$ 1
        };

        % --- Titre ---
        \node[anchor=center] at ([shift={(0cm,5cm)}]current page.center) {
            \fontsize{50}{60}\selectfont \textbf{GPIO et mesures de performances}
        };

        % --- Cadre de Note (Rouge) ---
        \node[anchor=north west] at ([shift={(1.2cm,-12.2cm)}]current page.north west) {
            \large Note :
        };

        % --- Logo (Bas Centre) ---
        % Remplace "logo-esisar" par le nom de ton fichier image
        \node[anchor=south] at ([shift={(0cm,3cm)}]current page.south) {
            \includegraphics[width=8cm]{figure/ESISAR-LOGO.png} 
        };
        
        \end{tikzpicture}
    \end{titlepage}

	\thispagestyle{fancy}
	

	%\section{Notions}
	
	%\begin{rappel}
		%\textbf{Loi d'Ohm :} $U = RI$
		
		%\textbf{Resistance de charge} : résistance qui limite le courant pendant le chargement/déchargement des condensateurs
		
		%\textbf{Circuit RC :} La constante de temps est donnée par $\tau = RC$
	%\end{rappel}
	
\newpage
\section{Préparation individuelle}

\begin{itemize}
    \item À partir de la documentation de la carte, identifier la référence du microcontrôleur ainsi que le port et le pin sur lequel est branché la LED.
    \begin{solution}
        \noindent
        \parbox[t]{6cm}{
            \textbf{Microcontrôleur :} ATmega328P \\
            \textbf{Port LED :} PB5 \\
            \textbf{Pin LED :} D13
        }
        \hfill
        \raisebox{-\height+\ht\strutbox}{
            \includegraphics[width=7.5cm]{figure/Port LED.png}
        }\\\\
        \textbf{Source :} Doc pinout.pdf Page 1
    \end{solution}

    \item Quelle est la différence entre le port PB5 et le port D13 ?
    \begin{solution}
        PB5 est le pin du bloc ATmega328P (adresse matérielle réelle), D13 celui de l'arduino (alias logique def par la carte).
    \end{solution}

    \item À partir de la documentation du microcontrôleur (chapitre 13), identifier les registres à configurer et les valeurs associées pour avoir le pin de la LED en sortie haute.
    \begin{solution}
        Pour mettre la broche PB5 à l'état haut, il faut configurer 2 bits :\\
        \begin{enumerate}
            \item DDB5 (bit 5 du registre DDRB) --> mettre à 1
                \begin{lstlisting}[language=C++] 
DDRB |= (1 << DDB5); // PB5 en sortie
                \end{lstlisting}\\
            \item PORTB5 (bit 5 du port PORTB) --> mettre à 1
                \begin{lstlisting}[language=C++] 
PORTB |= (1 << PORTB5); // PB5 a l'etat haut
                \end{lstlisting}\\
        \end{enumerate}
        Code arduino : 
        \begin{lstlisting}[language=c++]
pinMode(13, OUTPUT); digitalWrite(13, HIGH);
        \end{lstlisting}\\\\
        \textbf{Source :} Doc ATmega328P\_Datasheet.pdf Section 13.2.1 \& 13.4.2 
    \end{solution}

    \item À partir de la documentation assembleur et de la documentation microcontrôleur, comment utiliser l'instruction SBI pour inverser le port PB5 ? Donner la ligne assembleur correspondante.
    \begin{solution}
        D'après la datasheet ATmega328P (section 13.2.2), écrire 1 dans le bit du registre PINB fait basculer physiquement la sortie, quel que soit son état actuel.
        \begin{lstlisting}[language=asm]
SBI 0x03, 5
        \end{lstlisting}\\\\
        \textbf{Source :} Doc assembleur Section 99 \& Doc ATmega328P\_Datasheet.pdf Section 13.2.2
    \end{solution}

    \item À partir de la documentation de la carte, identifier la référence de l'oscillateur cristal.
    \begin{solution}
        Il s'agit d'un ECS-160-20-4X-DU Oscillator, correspondant à la ref Y1 de la carte.\\\\
        \textbf{Source :} Doc ArduinoUnoR3\_Datasheet.pdf Section 3.1
    \end{solution}

    \item À partir de la documentation de l'oscillateur (à chercher vous même), identifier ce à quoi chacun des champs de la référence correspond. Quelle est la précision du cristal à 25\,\textdegree{}C au bout de 1 an ?
    \begin{solution}
        \begin{tabular}{|c|c|c|}
            \hline
            Champ & Valeur & Signification \\
            \hline
            \hline
            ECS & NA & Fabricant (ECS Inc.)\\
            \hline
            160 & 16.000 MHz & Code fréquence (160 --> 16.000 MHz)\\
            \hline
            20 & 20 pF & Capacité de charge \\
            \hline
            4X & 3,5 mm & Hauteur du boîtier \\
            \hline
            D & $\pm$ 100 ppm & Stabilité en fréquence \\
            \hline
            U & -55°C \~ +125°C & Plage de température de fonctionnement \\
            \hline
        \end{tabular}\\
        \textbf{Précision à 25°C au bout d'un an :} `Aging (First Year) @+25°C$\pm$3°C:$\pm$5ppm`\\
        C'est-à-dire que le cristal varie au maximum de $\pm$5ppm par rapport à sa fréquence nominale (16 MHz). Soit une dérive maximale de $\pm$80Hz.\\\\
        \textbf{Source :} Doc ECS-160-20-4X-DU\_Datasheet.pdf Section "Part Numbering Guide" & "Operation Conditions/Electrical Characteristics"
    \end{solution}
\end{itemize}

\section{GPIO}

Nous allons comparer la différence de performance entre 4 méthodes pour faire clignoter une LED : deux méthodes en C avec les registres matériels, une en assembleur avec les registres également et la dernière en C avec une librairie plus haut-niveau (plus abstraite).

\begin{itemize}
    \item Faire la configuration des registres pour obtenir le mode Push-Pull et allumer la LED. Pour cela, il faut inclure l'en-tête spécifique à notre Arduino UNO R3 pour avoir accès aux registres. Le compilateur avr-gcc fournit une interface unique \texttt{\#include <avr/io.h>} qui lui même inclut le bon en-tête de notre cible (ici \texttt{avr/iom328p.h}).
    \begin{solution}
        \begin{lstlisting}[language=C++]
    void setup(){
        DDRB |= (1 << DDB5);
        PORTB |= (1 << PORTB5);
    }
    void loop(){
        
    }
        \end{lstlisting}
    \end{solution}

    \item Dans une boucle, faire clignoter la LED en utilisant le registre d'écriture et mesurer la fréquence à l'oscilloscope.
    \begin{solution}
        \begin{lstlisting}[language=C++]
   void setup(){
        pinMode(13, OUTPUT);        
    }
    void loop(){
        digitalWrite(13, HIGH);
        delay(500);
        digitalWrite(13, LOW);
        delay(500);
    }
        \end{lstlisting}
    \end{solution}

    \item De même en utilisant cette fois le registre permettant d'inverser l'état du pin (toggle).
    \begin{solution}
            \begin{lstlisting}[language=C++]
   void setup(){
        DDRB |= (1 << DDB5);
        
    }
    void loop(){
        PORTB |= (1 << PORTB5);
        delay(500);
        PORTB &= (0 << PORTB5);
        delay(500);
    }
        \end{lstlisting}
    \end{solution}

    \item Laquelle des 2 méthodes est la plus rapide ? Pourquoi ? Vous pouvez regarder le code assembleur de votre binaire en utilisant l'utilitaire \texttt{avr-objdump}. L'IDE Arduino génère le binaire compilé en faisant \texttt{Sketch > Export compiled binary} dans le dossier de votre projet Arduino :
    \begin{verbatim}
avr-objdump -DS yourprogram.ino.elf > yourprogram.txt
    \end{verbatim}
    Le drapeau \texttt{-D} indique de désassembler toutes les sections (y compris les données) et le drapeau \texttt{-S} entremêle l'assembleur généré avec votre code source.
    \begin{solution}
        La deuxième méthode est la plus rapide car nous utilisons simplement un bit shift alors que avec \texttt{digitalWrite} nous faisons beaucoup plus d'instructions.
    \end{solution}

    \item De même en utilisant les instructions \texttt{sbi} et \texttt{rjmp} dans un bout de code assembleur inline. La syntaxe est :
    \begin{verbatim}
asm volatile("ligne asm1\n"
             "ligne asm2\n"
             "ligne asm3\n");
    \end{verbatim}
    \begin{solution}
        \begin{lstlisting}[language=C++]
void setup(){
    pinMode(13, OUTPUT);
}
void loop(){
    asmvolatile(
        "sbi 0x03, 5 \n\t"
        "rjmp .-4 \n\t"
    );
}
        \end{lstlisting}
    \end{solution}

    \item Faire clignoter la LED en utilisant la librairie Arduino et la fonction \texttt{digitalWrite} et mesurer la fréquence. En sachant que le microcontrôleur fonctionne à 16\,MHz, combien de cycles d'horloge sont nécessaires pour un simple toggle ? Comparer avec les solutions précédentes.
    \begin{solution}
    Avec la fonction \texttt{digitalWrite}, nous mesurons une fréquence de 150\,kHz et de 2\,MHz avec l'accès direct des registres (fréquence max de 4\,MHz).

    \begin{itemize}
        \item Période du signal : \( T = \frac{1}{f} \)
        \item Un toggle = une demi-période : \( T_{\text{toggle}} = \frac{T}{2} \)
    \end{itemize}

    \medskip
    \textbf{Calcul des cycles par toggle :}

    \begin{itemize}
        \item \texttt{digitalWrite} : \( f = 150\,\text{kHz} \)
        \[
        N = \frac{T_{\text{toggle}}}{T_{\text{CPU}}}
        = \frac{1}{2 f \cdot f_{\text{CPU}}}
        = \frac{1}{2 \times 150\times10^{3} \times 16\times10^{6}}
        \approx \boxed{208 \text{ cycles}}
        \]

        \item Accès direct registre : \( f = 2\,\text{MHz} \)
        \[
        N = \frac{1}{2 \times 2\times10^{6} \times 16\times10^{6}}
        = \boxed{4 \text{ cycles}}
        \]

        \item Fréquence max : \( f = 4\,\text{MHz} \)
        \[
        N = \frac{1}{2 \times 4\times10^{6} \times 16\times10^{6}}
        = \boxed{2 \text{ cycles}}
        \]
    \end{itemize}

    \begin{center}
    \begin{tabular}{|l|c|c|}
        \hline
        \textbf{Méthode} & \textbf{Cycles / toggle} & \textbf{Fréquence} \\
        \hline
        \texttt{digitalWrite} & \(\sim 208\) & 150\,kHz \\
        Accès direct registre & 4 & 2\,MHz \\
        Assembleur & 4 & 2\,MHz \\
        Limite théorique & 2 & 4\,MHz \\
        \hline
    \end{tabular}
    \end{center}

    \medskip
    \textbf{Conclusion :} L'accès direct au registre est environ \( \frac{208}{4} \approx 52 \) fois plus rapide que \texttt{digitalWrite}, ce qui illustre le coût de l'abstraction Arduino.
\end{solution}

    \item Faire clignoter la LED en utilisant la librairie Arduino et une attente logicielle pour obtenir une fréquence de 1\,kHz. De manière aléatoire (par exemple 1 fois sur 2), ajouter une sortie sur le port série réglé à 9\,600 bauds d'un message quelconque (au moins 10 caractères). Quel est l'effet sur la fréquence de clignotement de la LED ?
    \begin{solution}
            \begin{lstlisting}[language=C++]
   void setup(){
        DDRB |= (1 << DDB5);
        Serial.begin(9600);
        
    }
    void loop(){
        PORTB |= (1 << PORTB5);
        delayMicroseconds(500);
        PORTB &= (0 << PORTB5);
        delayMicroseconds(500);

        //Serial.println("abcdefghijk");
    }
        \end{lstlisting}
    Avec l'affichage d'une chaine de charactère fait passer la fréquence de 1kHz à 70Hz.
    \end{solution}

    \item Même chose que précédemment mais sans attente logicielle.
    \begin{solution}
        Sans les attente logicielle, nous avons toujours une fréquence de 70Hz.
    \end{solution}

    \item Comparer les 2 codes précédents. Quels sont les avantages et inconvénients de ces 2 codes ?
    \begin{solution}
        \texttt{millis()} n'interrompt pas le code et permet de faire plusieurs tâches en parallèle. Le \texttt{delay()} interrompt tout le code.
    \end{solution}
\end{itemize}

\section{Préparation individuelle}

\begin{itemize}
    \item À partir de la documentation de la librairie de la fonction \texttt{attachInterrupt}, quels pins sont disponibles pour des interruptions ? Identifier les pins correspondants du microcontrôleur.
    \begin{solution}
    les interruptions externes de l'arduino sont sur les pins numériques D2 et D3.
    
    \end{solution}

    \item À partir de la documentation du microcontrôleur, quelles sont les interruptions disponibles sur les pins précédemment identifiés ? Sur quel pin se trouve l'interruption externe 0 ?

    
    \textit{Indice : regarder les fonctions alternatives des GPIO.}
    \begin{solution}
    
    \centering 
    \begin{tabular}{|c|c|c|c|} \hline 
    \textbf{Pin Arduino} & \textbf{Pin ATmega328P} & \textbf{INT externe} & \textbf{Alt pinout} \\ \hline 
    D2 & PD2 & INT0 & PCINT18 \\ \hline 
    D3 & PD3 & INT1 & PCINT19 \\ \hline 
    \end{tabular} \\ 
    Lien entre Arduino et les interruptions de l'ATmega328P.
    
    \end{solution}

    \item Quelle est la différence entre les interruptions INT0 et PCINT0 ?
    \begin{solution}
    INT0 est une interruption externe dédiée. Elle est associée à la broche PD2 et dispose de son propre vecteur d'interruption. Son mode de déclenchement peut être configuré afin de détecter : 
    \begin{itemize} 
        \item un niveau bas ; 
        \item un changement d'état logique ; 
        \item un front descendant ; 
        \item un front montant.
    \end{itemize} 
    Le choix du mode de déclenchement est réalisé à l'aide des bits \texttt{ISC01} et \texttt{ISC00} dans le registre \texttt{EICRA}. 
    \texttt{PCINT0} est une interruption de type \textit{Pin Change Interrupt}. Elle est située sur la pin PB0, correspondant au pin numérique D8 de l'Arduino Uno. 
    Les interruptions \textit{Pin Change} sont regroupées. Ainsi, \texttt{PCINT0} partage son vecteur d'interruption avec les autres interruptions du même groupe (8 interruptions). Si une interruption est toggle, un des trois pin \texttt{PCIx} correspondant est set à 1.
    \begin{itemize}
        \item \texttt{PCICR} : activer/désactiver le groupe d'interruptions
        \item \texttt{PCIFR} : Flag par groupement d'interruptions
        \item \texttt{PCMSKx} : activer/désactiver unitairement les interruptions 
    \end{itemize}
    
    Elles permettent de détecter un changement d'état logique, mais ne permettent pas directement de sélectionner uniquement un front montant ou uniquement un front descendant.
    \\
    
    \centering 
    \begin{tabular}{|l|c|c|} \hline
        & \textbf{INT0} & \textbf{PCINT0} \\ \hline
        Broche & PD2 / D2 & PB0 / D8 \\ \hline
        Type & Interruption dédiée & Pin Change Interrupt \\ \hline
        Déclenchement & Niveau ou front configurable & Changement logique \\ \hline
        Vecteur & Vecteur dédié & Vecteur partagé \\ \hline
        Registres principaux & EICRA, EIMSK & PCICR, PCMSK0 \\ \hline
    \end{tabular} \\
    Comparaison entre INT0 et PCINT0.
        
    \end{solution}

    \item À partir du chapitre sur les interruptions externes dans la documentation microcontrôleur, identifier les registres et les valeurs requises pour configurer l'interruption externe 0 (INT0) sur front montant.
    \begin{solution}
    Dans le registre \texttt{EICRA} les bits \texttt{ISC01} \&  \texttt{ISC00} en rising edge donc 1 et 1 respectivement. Il faut aussi penser à configurer tous les autres registres nécessaire à l'interruption souhaitée \texttt{EIMSK}...\\

    \centering
    \begin{tabular}{|c|c|l|} \hline 
    \textbf{ISC01} & \textbf{ISC00} & \textbf{Déclenchement} \\ \hline
        0 & 0 & Niveau bas \\ \hline 
        0 & 1 & Changement logique \\ \hline
        1 & 0 & Front descendant \\ \hline
        1 & 1 & Front montant \\ \hline
    \end{tabular} \\
    Configuration du déclenchement de l'interruption INT0.
    
    \end{solution}

    \item À partir de la documentation du compilateur avr-gcc, identifier les noms des vecteurs d'interruption pour les interruptions INT0 et PCINT2.
    \begin{solution}
    
    \begin{itemize}
        \item \textbf{INT0}  le nom du vecteur est : \texttt{INT0\_vect}
        \item \textbf{PCINT2} le nom du vecteur est : \texttt{PCINT2\_vect}
    \end{itemize}
    
    Les fonction d'interruption peuvent être écrite sous la forme : 
    \begin{verbatim} 
    ISR(Nom_Vecteur) 
    {
    // Traitement de l'interruption INT0 
    } 
    \end{verbatim}
    
    \end{solution}
\end{itemize}

\section{Interruption}

\begin{itemize}
    \item Configurer l'interruption externe 0 sur front montant.
    \begin{solution}
        \begin{lstlisting}[language=C++]
void setup(){
    pinMode(13, OUTPUT);
    EICRA |= (1 << ISC00);  //front
    EICRA |= (1 << ISC01);  //montant
    EIMSK |= (1 << INT0);   //Enable
    EIFR |= (1 << INTF0);   //flag
}
void loop(){

}    
        \end{lstlisting}
    \end{solution}

    \item Écrire la fonction d'interruption telle que : \texttt{ISR([vector\_name]) \{...\}}, en remplaçant \texttt{[vector\_name]} par le nom identifié dans la préparation, et allumer la LED dans cette fonction. Éteindre la LED depuis la fonction \texttt{loop}.
    \begin{solution}
        \begin{lstlisting}[language=C++]
#include <avr/io.h>
#include <avr/interrupt.h>

ISR(INT0_vect) {
    PORTB |= (1 << PORTB5); 
}

void setup() {
    pinMode(13, OUTPUT);
    PORTB |= (1 << PORTB5); 
    EICRA |= (1 << ISC01) 
    EICRA |= (1 << ISC00);
    EIMSK |= (1 << INT0); 
    EIFR  |= (1 << INTF0);
    sei();
}

void loop() {
    PORTB &= (0 << PORTB5); 
}
        \end{lstlisting}
    \end{solution}

    \item Mesurer le temps de réponse à l'oscilloscope en faisant l'écriture du registre et en passant par la fonction \texttt{digitalWrite}.
    \begin{solution}
        Le temps de réponse de l'oscilloscope est de 1,4$\micro$s.\\
        Ce qui correspond à 22 cycles.\\
        Avec \texttt{digitalWrite}, le temsp de réponse est de 4,1.\\
        Ce qui correspond à 65 cycle.
    \end{solution}

    \item Commenter votre code précédent et faire l'interruption en utilisant la fonction \texttt{attachInterrupt}.
    \begin{solution}
        \begin{lstlisting}[language=C++]
void ISR(){    
    PORTB &= (0 << PORTB5);
}
void setup(){
    pinMode(13, OUTPUT);
    EICRA |= (1 << ISC00);  //front
    EICRA |= (1 << ISC01);  //montant
    EIMSK |= (1 << INT0);   //Enable
    EIFR |= (1 << INTF0);   //flag
    PORTB &= (0 << PORTB5);
    attachInterrupt(digitalPinToInterrupt(2), ISR, RISING);
    sei();
}
void loop(){
    PORTB |= (1 << PORTB5);
}    
        \end{lstlisting}
    \end{solution}

    \item Mesurer le temps de réponse à l'oscilloscope. Conclure.
    \begin{solution}
        Le temps de réponse de l'oscilloscope est 3,1$\micro$s. Ce qui correspond à 50 cycles.\\\\
        Ainsi, le \texttt{attachInterrupt} utilise 23 cycles de plus.
    \end{solution}
\end{itemize}

\section{Timers}

\begin{itemize}
    \item Quelle est la taille du timer 1 ? Quelle est la division d'horloge maximale du prescaler ? En déduire la période maximale du timer 1.
    \begin{solution}
        Le timer 1 est un timer 16 bits (il peut compter de 0 à 65 535). Il possède plusieurs valeurs de prescaler : 1, 8, 64, 256 et 1024. La plus grande est donc 1024.\\
        $P_{ériode_{max}} = 1024 * \frac{65535}{16*10^{6}}=4,19 s$
    \end{solution}

    \item En utilisant le mode normal du timer 1, quelle doit-être la division d'horloge pour être le plus proche d'une base de temps de 1 seconde ? Implémenter le mode normal et mesurer la fréquence à l'oscilloscope. La fonction d'interruption se déclare comme suit : \texttt{ISR(TIMER1\_OVF\_vect) \{...\}}
    \begin{solution}
        $P_{ériode_{1s}} = 256*\frac{65535}{16*10^{6}} = 1,05 s$
        Le prescaler qu'il faut est 256 pour avoir un basse de temps de environ 1 seconde.
        \begin{lstlisting}[language=C++]
#include <avr/io.h>
#include <avr/interrupt.h>

volatile uint16_t compteur = 0;

void setup() {
    TCCR1A = 0;
    TCCR1B = 0;
    TCCR1B |= (1 << CS12);
    TIMSK1 |= (1 << TOIE1);
    sei();
}
ISR(TIMER1_OVF_vect) {
    PORTB ^= (1 << PORTB5);
}
        \end{lstlisting}
        Il y a un débordement car notre période n'est pas exactement de 1s.
    \end{solution}

    \begin{solution}
        En mode CTC (\textit{Clear Timer on Compare Match}), le compteur \texttt{TCNT1}
        est remis à zéro lorsqu'il atteint la valeur contenue dans le registre
        \texttt{OCR1A}. La période entre deux comparaisons vaut donc :
        \[
            T_{\mathrm{CTC}}
            = \frac{N\left(\mathrm{OCR1A}+1\right)}{f_{\mathrm{CPU}}},
        \]
        où \(N\) est le facteur de division du prescaler et
        \(f_{\mathrm{CPU}}=16\,\mathrm{MHz}\).

        On souhaite obtenir une base de temps de \(1\,\mathrm{s}\), donc :
        \[
            1
            = \frac{N\left(\mathrm{OCR1A}+1\right)}{16\times 10^6}.
        \]
        D'où :
        \[
            \mathrm{OCR1A}+1=\frac{16\times 10^6}{N}.
        \]

        Avec un prescaler de \(256\) :
        \[
            \mathrm{OCR1A}+1
            = \frac{16\times 10^6}{256}
            = 62500,
        \]
        et donc :
        \[
            \boxed{\mathrm{OCR1A}=62499}.
        \]

        Cette valeur est inférieure à \(65535\), donc elle est compatible avec
        le Timer~1 qui est un timer 16 bits. On obtient alors exactement :
        \[
            T_{\mathrm{CTC}}
            = \frac{256\times 62500}{16\times 10^6}
            = \boxed{1\,\mathrm{s}}.
        \]

        Le registre qui contient la valeur maximale en mode CTC 4 est donc
        \boxed{\texttt{OCR1A}}.

        Pour sélectionner le mode CTC 4, il faut configurer
        \(WGM13:0=0100\), donc mettre \texttt{WGM12} à 1 et laisser les autres
        bits WGM à 0.
        
        \end{solution}
        
        \begin{solution}
        \begin{lstlisting}[language=C++]
void setup() {
    // LED D13 = PB5 en sortie
    DDRB |= (1 << DDB5);

    // Timer1 : mode CTC (mode 4)
    TCCR1A = 0;
    TCCR1B = 0;
    TCCR1B |= (1 << WGM12);

    OCR1A = 62499; //  interruption toutes les 1 s
    // Prescaler = 256
    TCCR1B |= (1 << CS12);
    // Autorisation de l'interruption Compare Match A
    TIMSK1 |= (1 << OCIE1A);
    sei();
}

ISR(TIMER1_COMPA_vect) {
    PINB |= (1 << PINB5);
}

void loop() {}
        \end{lstlisting}

        L'interruption est déclenchée toutes les \(1\,\mathrm{s}\).
        Comme la LED est inversée à chaque interruption, il faut deux
        interruptions pour obtenir une période complète du signal de sortie :
        \[
            T_{\mathrm{LED}}=2\,\mathrm{s},
            \qquad
            f_{\mathrm{LED}}=\boxed{0{,}5\,\mathrm{Hz}}.
        \]
        L'oscilloscope doit donc mesurer environ \(0{,}5\,\mathrm{Hz}\) sur la
        LED, tandis que les évènements de comparaison ont lieu à \(1\,\mathrm{Hz}\).

        \medskip
        \textbf{Remarque :} si l'on souhaite directement obtenir un signal carré
        de \(1\,\mathrm{Hz}\) en inversant la sortie à chaque interruption, il
        faut provoquer une interruption toutes les \(0{,}5\,\mathrm{s}\), soit
        \(\mathrm{OCR1A}=31249\) avec le même prescaler de 256.
    \end{solution}

    \item Faire de même sans interruption et en utilisant l'un des modes PWM.
    \begin{solution}
        L'objectif est maintenant de générer directement un signal périodique
        d'environ \(1\,\mathrm{Hz}\) sans utiliser d'interruption. Pour cela,
        on peut exploiter une sortie matérielle du Timer~1.

        \medskip
        On choisit le mode \textbf{Fast PWM 14}, pour lequel :
        \[
            WGM13:0 = 1110
        \]
        et la valeur maximale du compteur (\textit{TOP}) est contenue dans
        \texttt{ICR1}.

        Dans ce mode, la fréquence PWM vaut :
        \[
            f_{\mathrm{PWM}}
            =
            \frac{f_{\mathrm{CPU}}}
                 {N\left(1+\mathrm{ICR1}\right)}.
        \]

        Avec \(f_{\mathrm{CPU}}=16\,\mathrm{MHz}\), un prescaler
        \(N=256\) et une fréquence souhaitée de \(1\,\mathrm{Hz}\), on obtient :
        \[
            1
            =
            \frac{16\times10^6}
                 {256\left(1+\mathrm{ICR1}\right)}.
        \]
        Donc :
        \[
            1+\mathrm{ICR1}
            =
            \frac{16\times10^6}{256}
            =
            62500,
        \]
        soit :
        \[
            \boxed{\mathrm{ICR1}=62499}.
        \]

        Pour obtenir un rapport cyclique de \(50\,\%\), on choisit :
        \[
            \mathrm{OCR1A}+1
            =
            \frac{\mathrm{ICR1}+1}{2}
            =
            31250,
        \]
        donc :
        \[
            \boxed{\mathrm{OCR1A}=31249}.
        \]

        \textbf{Attention :} la LED intégrée de l'Arduino Uno est connectée à
        PB5 (D13), mais PB5 n'est pas une sortie PWM matérielle du Timer~1.
        La sortie \texttt{OC1A} du Timer~1 est située sur PB1, c'est-à-dire sur
        la broche Arduino \textbf{D9}. Pour observer le signal sans
        interruption, il faut donc mesurer sur D9 (ou utiliser OC1B sur D10).
        \end{solution}
        \begin{solution}
        
        \begin{lstlisting}[language=C++]
void setup() {
    // D9 = PB1 = OC1A en sortie
    DDRB |= (1 << DDB1);

    // Remise a zero de la configuration du Timer1
    TCCR1A = 0;
    TCCR1B = 0;

    // PWM mode 14 :
    TCCR1A |= (1 << WGM11);
    TCCR1B |= (1 << WGM13) | (1 << WGM12);

    // Sortie OC1A non-inversee :
    TCCR1A |= (1 << COM1A1);
    ICR1 = 62499;
    // Rapport cyclique = 50 %
    OCR1A = 31249;
    // Prescaler = 256
    TCCR1B |= (1 << CS12);
}

void loop() {}
        \end{lstlisting}

        La fréquence obtenue est :
        \[
            f_{\mathrm{PWM}}
            =
            \frac{16\times10^6}
                 {256\times62500}
            =
            \boxed{1\,\mathrm{Hz}}.
        \]

        La période est donc :
        \[
            \boxed{T=1\,\mathrm{s}},
        \]
        avec un rapport cyclique d'environ \(50\,\%\).

        \medskip
        \textbf{Conclusion :} cette solution est plus efficace que la solution
        avec interruption car le signal est entièrement généré par le matériel
        du Timer~1. Le processeur n'a pas besoin d'exécuter une ISR à chaque
        période et reste disponible pour d'autres traitements.
    \end{solution}
\end{itemize}
    
\end{document}
